% Standalone mathematical fragment. Requires amsmath, amssymb, amsthm,
% and theorem, lemma, proposition, corollary, definition, remark environments.
% The parent paper should credit the two-row package to the incoming
% ipra_variable_substitution/paper/variable_substitution_proof.tex and the
% compound-envelope repair to Research/full_normalization/report.tex.
% No proof-assistant verification is claimed by this fragment.

\section{Equality at the discharge of a numerical witness}
\label{fus:section}

The change is a uniform strengthening of paired elimination. Equality
substitution is performed only while discharging the fresh numerical witness
of that elimination. An arbitrary equality inference elsewhere in a proof is
recovered by first giving its source a fresh name. The naming construction is
inherited from the preceding variable-substitution proposal; its final
substitution and witness discharge become one primitive inference here.

\subsection{Syntax, normal forms, and the four-pattern presentation}

Write $N$ for full computation normalization of finite scalar-iterator terms
\[
 t::=x\mid 0\mid S(t)\mid I(n,a;u.h),
\]
under the two context-closed reductions
\[
 I(0,a;u.h)\longrightarrow a,
 \qquad
 I(Sn,a;u.h)\longrightarrow h[I(n,a;u.h)/u].
\]
Only $u$ in $h$ is bound. All substitutions avoid capture. Numerical terms
of the calculus are the normal forms; every substituted instance is normalized.
In particular, $I(n,a;u.h)$ in a displayed formula denotes the normal form of
that constructor expression. Formulas and reusable clauses are
\[
 F::=(s=t)\mid D_{x,y}F,
 \qquad [A;y.B(y)],\qquad y\notin\operatorname{FV}(A).
\]
Normalization acts on numerical subterms and does not delete formula binders.
A context is a finite list of clause occurrences. Put $H=(0=0)$, and abbreviate
an ordinary formula assumption $F$ by the constant clause $[H;v.F]$, with $v$
unused. Context membership is by occurrence, so equal-valued clauses may still
have distinct positions.

The presentation $\mathcal F_4$ has the following four inference patterns.
\[
 \frac{}{\Gamma\vdash t=t}.\tag{R}\label{fus:refl}
\]
\[
 \frac{\Gamma\vdash A}{\Gamma\vdash N(B(t))},
 \qquad [A;y.B(y)]\in\Gamma.\tag{U}\label{fus:use}
\]
\[
 \frac{\Gamma,B(t,\eta)\vdash B(t,Sz)}
      {\Gamma\vdash D_{x,y}B(x,y)},\qquad
 \eta=\begin{cases}
 z,& B\text{ has equation root and }t\equiv0,\\
 0,&\text{otherwise}.
 \end{cases}\tag{I}\label{fus:intro}
\]
Here $z$ is fresh for $\Gamma$, $t$, and all passive matrix parameters. All
instances in this display are canonical; $t\equiv0$ is literal equality of
canonical syntax. Finally,
\[
\frac{
 \Gamma\vdash D_{x,y}B(x,y)
 \quad
 \Gamma,\mathcal B_b\vdash b=t
 \quad
 \Gamma,\mathcal B_b\vdash F
}{
 \Gamma\vdash N(F[t/b])
},
\qquad
 \mathcal B_b=[B(b,0);y.B(b,y)].
\tag{$E_{=}$}\label{fus:elim}
\]
The parameter $b$ is fresh for $\Gamma$ and the passive parameters of $B$.
The two branch proofs have the \emph{same} extended context. The term $t$ and
formula $F$ are well scoped in that branch, and the exported formula
$N(F[t/b])$ must be free of $b$. Thus the exported formula is well scoped in
the outer numerical context. The substitution is simultaneous: an inserted
copy of $t$ is not traversed again.

In particular, $t$ is allowed to contain $b$. The special choice $t=b$ is
essential below. The condition concerns the exported normal form, not a
syntactic prohibition of $b$ in $t$. These are all the primitive inferences.
There is no additional equality inference, computation axiom, logical axiom,
successor axiom, or induction rule. The two computation equations are built
into canonical syntax and substitution; their normalizer is part of the
trusted presentation and the measured checker.

\subsection{Direct soundness of the fused inference}

The ordinary readings are
\[
 (D_{x,y}B)^\circ
 =\exists x\bigl(B^\circ(x,0)\to\forall y\,B^\circ(x,y)\bigr),
 \qquad
 [A;y.B]^\circ=A^\circ\to\forall y\,B^\circ(y).
\]
Equality has its ordinary meaning. For soundness in the inherited ordinary
scalar theory, rule I uses numerical coverage in its case mode and literal-
equation induction in its progress mode. R and U have their indicated usual
readings. The only changed inference is $E_{=}$.

\begin{lemma}[Soundness of witness discharge with substitution]
\label{fus:soundness}
Rule~\eqref{fus:elim} is sound in intuitionistic first-order equality logic
over the inherited scalar arithmetic, for arbitrary $F$, including compound
formulas, and also when $t$ contains $b$.
\end{lemma}
\begin{proof}
Eliminate the existential in the ordinary reading of the major premise at
fresh $b$. Its local assumption is precisely the reading of $\mathcal B_b$.
The two supplied branch derivations give $b=t$ and $F$. Ordinary equality
substitution gives $F[t/b]$, and the iterator computation equations identify
this formula with its normal form. By the rule's side condition that normal
form has no free $b$, so ordinary existential elimination discharges $b$ and
its clause.

This equality step does not require $b\notin\operatorname{FV}(t)$.
For example, equality $b=t(b)$ supports simultaneous substitution of $t(b)$
for $b$ in any fixed formula. The condition needed for existential elimination
is the absence of $b$ from the \emph{result}, which the rule requires explicitly.
\end{proof}

\subsection{What can be recovered before any equality reasoning}

\begin{lemma}[Ordinary paired elimination]
\label{fus:ordinary-elim}
The ordinary rule
\[
 \frac{\Gamma\vdash D_{x,y}B(x,y)\qquad
       \Gamma,\mathcal B_b\vdash A}
      {\Gamma\vdash A},\qquad b\notin\operatorname{FV}(\Gamma,A),
\]
is derivable in $\mathcal F_4$.
\end{lemma}
\begin{proof}
Take $t=b$ in $E_{=}$. Its equality branch is reflexivity $b=b$, and its
other branch is the given proof of $A$. Since $A$ has no free $b$,
$N(A[b/b])=A$ literally.
\end{proof}

Weakening, injective fresh renaming, and ordinary-assumption cut are
admissible by direct induction on derivations. Freshen the eigenparameters
of I and $E_{=}$ before inserting proofs. The substitution source in
$E_{=}$ is the parameter bound by that very inference; inserting additional
parameters outside its scope does not change that source. A cut of an
ordinary assumption replaces uses of its constant clause by the supplied
proof. These arguments require no arbitrary numerical substitution theorem
for derivations and no equality congruence.

For an unused $w$, abbreviate $q_yA(y)=D_{w,y}A(y)$.
\begin{lemma}[Closure and numerical cases without equality]
\label{fus:cases}
Using only R, U, I, and derived ordinary paired elimination, one obtains
\[
 \frac{\Gamma,A(0)\vdash A(Sz)}{\Gamma\vdash q_yA(y)},
 \qquad
 \frac{\Gamma\vdash q_yA(y)\qquad\Gamma\vdash A(0)}
      {\Gamma\vdash A(t)}.
\]
Consequently, from $\Gamma\vdash A(0)$ and $\Gamma\vdash A(Sz)$, with
$z$ fresh for the retained context and passive parameters, one obtains
$\Gamma\vdash A(t)$ for every well-scoped term $t$.
\end{lemma}
\begin{proof}
For the first rule introduce the unused witness coordinate at $1=S0$.
Rule I is then in case mode at every matrix root. For the second rule use
ordinary paired elimination: since $w$ is unused, its opened clause is
$[A(0);y.A(y)]$. Rule U with the supplied zero proof and target $t$ gives
the result. To get numerical cases, weaken the successor proof by the unused
assumption $A(0)$ and compose the two constructions.
\end{proof}

\subsection{Equality recovered by a two-row witness package}

The inherited normalizing selector and parity terms are
\[
 C(v,a,c)=N(I(v,a;u.c)),\quad u\notin\operatorname{FV}(c),
 \qquad
 P(y)=N(I(y,0;u.C(u,1,0))).
\]
Their endpoint computations are literal canonical-syntax identities:
\[
 C(0,a,c)\equiv a,\quad C(Sv,a,c)\equiv c,
 \qquad P(1)\equiv1,\quad P(2)\equiv0.
\]
No induction or range theorem for $P$ is needed.

\begin{theorem}[Uncrossed transport from one fused discharge]
\label{fus:uncrossed}
For every fixed numerical template $h(w)$, the inference
\[
 \frac{\Gamma\vdash s=t\qquad\Gamma\vdash N(h(s))=r}
      {\Gamma\vdash N(h(t))=r}
\]
is derivable in $\mathcal F_4$. Apart from ordinary paired eliminations
recovered by $t=b$, it uses one equality-bearing fused discharge.
\end{theorem}
\begin{proof}
Choose $x,y$ fresh for all passive data and define the equation matrix
\[
 B(x,y)=\bigl(
 C(y,0,C(P(y),x,N(h(x))))
 =C(y,0,C(P(y),t,r))\bigr).
\]
Introduce $D_{x,y}B(x,y)$ at witness $s$. Its successor instance is
\[
 C(P(Sz),s,N(h(s)))=C(P(Sz),t,r).
\]
Apply Lemma~\ref{fus:cases} to the formula
\[
 A(v)=\bigl(C(v,s,N(h(s)))=C(v,t,r)\bigr)
\]
at target $P(Sz)$. Its zero instance is the supplied $s=t$, and every
successor instance is the supplied $N(h(s))=r$. Choose the case
eigenparameter fresh for the current local context, independently of $z$.
This derivation ignores the temporary assumption of I. It therefore works
whether $s$ is literal zero and I uses progress, or I uses cases.

Now apply $E_{=}$ to this package with fresh parameter $b$. The opened
clause is $[H;y.B(b,y)]$, since $B(b,0)\equiv H$. By U at $1$ and $2$,
using reflexivity for $H$, obtain the literal normalized instances
\[
 B(b,1)\equiv (N(h(b))=r),\qquad B(b,2)\equiv (b=t).
\]
Use the latter as the equality branch of $E_{=}$, and use
$F=(N(h(b))=r)$ as its formula branch. Normalization--substitution coherence
identifies the exported formula as
\[
 N\bigl((N(h(b))=r)[t/b]\bigr)\equiv(N(h(t))=r).
\]
The fresh $b$ does not occur in $t,r,h$, or the outer context, so the
exported formula is outer scoped. This identity is a metasyntactic
coherence identity, not an invocation of an unproved equality rule.
\end{proof}

\begin{corollary}[Symmetry, crossed transport, and ordinary equality]
\label{fus:equality}
The presentation derives symmetry, transitivity, term congruence, crossed
transport $s=t,\,h(s)=r\Rightarrow r=h(t)$, and transport in every formula
template.
\end{corollary}
\begin{proof}
In Theorem~\ref{fus:uncrossed}, take $h(w)=w$ and $r=s$; the premises
$s=t$ and reflexivity $s=s$ give $t=s$. Thus symmetry is obtained without
assuming it in the naming package. Uncrossed transport followed by symmetry
gives crossed transport. The standard two applications of crossed transport
give replacement in both sides of an equation; symmetry and root joining
give transitivity, and a reflexive minor gives congruence.

For transport through a compound template $D_{x,y}B(v,x,y)$, argue by
external induction on formula depth. From $s=t$ and $D_{x,y}B(s,x,y)$,
use ordinary paired elimination at a fresh $c$. Introduce
$D_{x,y}B(t,x,y)$ at witness $c$, which is not literal zero. Under its
local assumption $B(t,c,0)$, the induction hypothesis and symmetry give
$B(s,c,0)$. Specialize the old opened clause at $Sz$, obtaining
$B(s,c,Sz)$, and transport forward by the induction hypothesis to
$B(t,c,Sz)$. Close I and ordinary elimination. Every recursive transport
uses a strictly smaller formula. This argument does not require induction
on quantified formulas.
\end{proof}

\subsection{Repair of atomic case introduction at the zero witness}
\label{fus:repair}

Full computation normalization removes closed nonliteral aliases for zero.
The inherited compound-envelope construction repairs the missing case
without such an alias; it remains available in the new presentation before
using Theorem~\ref{fus:uncrossed}.

For an equation $E=(s=t)$, put
\[
 G_vE=(C(v,0,s)=C(v,0,t)),\qquad L(E)=q_vG_vE.
\]
The gate has literal endpoints $G_0E\equiv H$ and $G_{Sw}E\equiv E$.
Lemma~\ref{fus:cases} gives $E\vdash L(E)$ by closure introduction,
and $L(E)\vdash E$ by closure specialization at $1$ and reflexivity of
$H$. These are uniform derivations under any retained context.

Suppose $E(x,y)$ is an equation and a supplied derivation proves
$\Gamma,E(0,0)\vdash E(0,Sz)$. Set $A(x,y)=L(E(x,y))$.
Under $A(0,0)$, unwrap the envelope, insert the supplied derivation by
ordinary-assumption cut, and wrap the resulting equation. This proves
$A(0,Sz)$. Rule I at witness zero now uses cases because $A$ is compound,
so it gives $D_{x,y}A(x,y)$.

Open this formula by ordinary elimination at a fresh $c$. To introduce
$D_{x,y}E(x,y)$ at witness $c$, assume $E(c,0)$, wrap it as $A(c,0)$,
use the opened clause at $Sz$, and unwrap $A(c,Sz)$ to $E(c,Sz)$.
The witness $c$ is not literal zero, so I again uses cases. Closing both
inferences yields $D_{x,y}E(x,y)$. Hence unrestricted case introduction
is a derived rule. This is a separate completeness repair, not an extra
primitive induction principle.

\subsection{Exact sequent equivalence and the interpretation boundary}

Let $\mathcal T_5$ be the same canonical Eq/$D$ clause calculus with R, U,
I, ordinary paired elimination, and crossed term transport as its five
patterns. This is the full-normalization clausal control.

\begin{theorem}[Exact identity-language equivalence]
\label{fus:equivalence}
For every finite canonical clause context $\Gamma$ and canonical formula $F$,
\[
 \Gamma\vdash_{\mathcal F_4}F
 \quad\Longleftrightarrow\quad
 \Gamma\vdash_{\mathcal T_5}F.
\]
Both directions are effective transformations of finite proof trees and
leave the entire end sequent literally unchanged.
\end{theorem}
\begin{proof}
Translate a $\mathcal T_5$ proof recursively. Replace ordinary elimination
by Lemma~\ref{fus:ordinary-elim}, and replace each crossed transport by
Theorem~\ref{fus:uncrossed} followed by the derived symmetry construction.
R, U, and I are unchanged. The preliminary numerical-case constructions
used by the naming package were derived without either equality transport
or a translated transport node, so the recursive translation is well founded.

Conversely, the control derives whole-formula equality transport by the
same formula-depth argument as Corollary~\ref{fus:equality}. In a translated
$E_{=}$ branch, apply that transport to $b=t$ and $F$, obtaining
$N(F[t/b])$. Its stipulated $b$-freedom permits ordinary paired elimination.
All other patterns are shared. Freshen the finitely many local parameters
before each replacement.
\end{proof}

Composing Theorem~\ref{fus:equivalence} with the attributed canonical/raw
comparison and the scalar/iPRA arithmetic bridge gives full first-order
intuitionistic PRA, not merely its equations or closed sentences. Numerical
interpretations use the whole numerical domains and actual equality. The
new comparison is the identity on terms and formulas, so it preserves the
inherited definable numerical composite isomorphisms with graph $x=y$ and
both inherited formula round trips. Equiconsistency follows in particular
because the common designated contradiction $0=S0$ is unchanged.

\paragraph{Normalization is metasyntactic.}
The comparison uses the established termination, confluence, and
substitution coherence of the iterator reductions. Canonical syntax
recognition itself is a structural check. Normalizing a supplied substituted
instance is a separate trusted computation and must be included in the
complete checker size. There is no numerical variable ranging over function
descriptions and no object-language uniform evaluator or totality axiom.
Uniform normalization over arbitrary primitive-recursive descriptions has
no primitive-recursive bound in general; the standard diagonal obstruction
to a universal primitive-recursive evaluator applies. The proof translations
above are claimed to be effective, not automatically primitive recursive.
Explicit computation traces would give a different certificate interface
and would need to be counted separately.

\paragraph{Substitution of derivations.}
A primitive-tag-preserving numerical-substitution theorem is not assumed:
normalization can turn a witnessing term into literal zero and thereby
change I's selected mode. The envelope repair or transfer through the
substitution-stable canonical reference calculus supplies the corresponding
admissibility result. The substitution inside $E_{=}$ is an operation on
its proved formula, justified by equality, not an attempt to substitute
blindly into the branch derivation.

\subsection{Three instruction patterns with an explicit fixed axiom clause}
\label{fus:three}

There is also a uniform-clause presentation $\mathcal F_3$. Its contexts
contain clauses $[\vec A;y.B(y)]$ with zero or one antecedent. Replace U by
\[
 \frac{\{\,\Gamma\vdash A:A\text{ is the antecedent occurrence in }\vec A\,\}}
      {\Gamma\vdash B(t)}.
\tag{$U_{0,1}$}\label{fus:uniform-use}
\]
For an empty antecedent list this inference has no premise. Every proof
starts with the one explicit fixed axiom clause
\[
             \mathcal R=[\,;y.y=y].\tag{Ax}\label{fus:axiom}
\]
Ordinary local assumptions $F$ are now unconditional constant clauses
$[\,;v.F]$. The I pattern is otherwise unchanged. The $E_{=}$ pattern opens
the singleton-premise clause $[[B(b,0)];y.B(b,y)]$ and has exactly the same
two branch conclusions, substitution, and scope condition. There is no R
instruction: $t=t$ is a use of $\mathcal R$ at $t$.

\begin{theorem}[Uniform clauses preserve all consequences]
\label{fus:uniform-equivalence}
The three instruction patterns $U_{0,1}$, I, $E_{=}$, together with the fixed
axiom clause $\mathcal R$, give the same theory as $\mathcal F_4$. The
comparison covers arbitrary retained clauses with zero or one antecedent,
not only clauses reachable in a closed proof.
\end{theorem}
\begin{proof}
For a context $\Delta$ of optional-premise clauses, define $L(\Delta)$
occurrence by occurrence by
\[
 [\,;y.B]\longmapsto[H;y.B],\qquad
 [[A];y.B]\longmapsto[A;y.B].
\]
Exclude the distinguished initial axiom occurrence $\mathcal R$ from this
translation; its uses are translated to R. Then for every $F$,
\[
 \mathcal R,\Delta\vdash_{\mathcal F_3}F
 \quad\Longleftrightarrow\quad
 L(\Delta)\vdash_{\mathcal F_4}F.
\]
Forwards, an unconditional U use becomes an ordinary U use supplied with
R at zero. A singleton U use is unchanged. The clauses opened by $E_{=}$
correspond exactly, as do I's local assumptions after this translation.
An axiom-clause use becomes R at its selected term. Translate the two
branch proofs of each fused elimination recursively.

Backwards, R is an empty-premise use of $\mathcal R$. For a U use selecting
a clause that came from an unconditional occurrence, omit the supplied
proof of $H$; singleton occurrences keep their premise proof. Local
ordinary assumptions introduced by I correspond to unconditional clauses
in $\mathcal F_3$, so their U uses are handled in the same way.
The fused elimination and introduction have unchanged targets and
side conditions. Retain the origin of each clause occurrence during this
translation, including when different origins have the same clause value.

The same argument applies to matched controls that have separate equality
transport and ordinary paired elimination. Thus the fusion theorem holds
with the optional-premise context syntax as well. All end formulas remain
identical, and the ordinary meanings of corresponding contexts agree.
\end{proof}

The number three counts executable inference patterns after reflexivity
has been moved to one explicitly listed initial axiom clause. It does not
assert that reflexivity was proved without an axiom. The mathematical
reduction specific to this report is the localization of all equality
elimination to the discharge of a numerical witness. The optional-premise
clause interface and the initial reflexivity clause are inherited choices
of presentation. Neither rule counting nor a smaller matched checker is
a proof of an absolutely simplest axiomatization over all languages,
codings, or implementations.
